% !PRINTED 139
\weberpage{139}

Wir können uns hiernach auf die Betrachtung \gesperrt{homogener, symmetrischer Functionen} beschränken, aus denen alle anderen sich zusammensetzen lassen.

Hiernach ist es leicht, die allgemeine Form einer symmetrischen Function anzugeben. Wir erhalten sie, wenn wir in einem Gliede einer solchen Function
\[
 \alpha_1^{\nu_1}\alpha_2^{\nu_2}\cdots\alpha_n^{\nu_n}
\]
(ähnlich wie bei der Bildung der Determinanten) die unteren Indices auf alle mögliche Art permutiren und die Summe aller so gebildeten Glieder nehmen. Eine solche Function können wir einen \gesperrt{Elementarbestandtheil} einer symmetrischen Function nennen. Nehmen wir mehrere solche Elementarbestandtheile, multipliciren sie mit beliebigen, von den $\alpha$ unabhängigen Factoren und addiren sie, so erhalten wir die allgemeinste symmetrische Function. Die Anzahl der Glieder eines dieser Elementarbestandtheile ist, wenn die Exponenten $\nu_1,\nu_2,\ldots,\nu_n$ alle von einander verschieden sind, $\Pi(n)$; wenn aber ein und derselbe Exponent $\nu$ mehrmals vorkommt, so hat man die Permutationen, die keine verschiedenen Glieder geben, wegzulassen. Es ist z. B. bei drei Veränderlichen, wenn $\nu_1,\nu_2,\nu_3$ verschieden sind, ein Elementarbestandtheil:
\[
\begin{aligned}
 &\alpha_1^{\nu_1}\alpha_2^{\nu_2}\alpha_3^{\nu_3}
+\alpha_1^{\nu_1}\alpha_3^{\nu_2}\alpha_2^{\nu_3}
+\alpha_2^{\nu_1}\alpha_1^{\nu_2}\alpha_3^{\nu_3}
+\alpha_2^{\nu_1}\alpha_3^{\nu_2}\alpha_1^{\nu_3}\\
 &\qquad+\alpha_3^{\nu_1}\alpha_1^{\nu_2}\alpha_2^{\nu_3}
+\alpha_3^{\nu_1}\alpha_2^{\nu_2}\alpha_1^{\nu_3},
\end{aligned}
\]
wenn aber $\nu_3=\nu_2$ ist,
\[
 \alpha_1^{\nu_1}\alpha_2^{\nu_2}\alpha_3^{\nu_2}
+\alpha_2^{\nu_1}\alpha_1^{\nu_2}\alpha_3^{\nu_2}
+\alpha_3^{\nu_1}\alpha_1^{\nu_2}\alpha_2^{\nu_2}.
\]

Das einfachste Beispiel einer symmetrischen Function ist die Summe der Variablen
\[
 \alpha_1+\alpha_2+\alpha_3+\cdots+\alpha_n.
\]
Ebenso gehört das Product $\alpha_1\alpha_2\cdots\alpha_n$ dazu.

Diese beiden sind die extremen Fälle einer Reihe von symmetrischen Functionen, die wir die \gesperrt{symmetrischen Grundfunctionen} nennen und folgendermaassen erhalten.

Das Product
\begin{equation}
 f(x)=(x-\alpha_1)(x-\alpha_2)\cdots(x-\alpha_n)
\tag{1}
\end{equation}
ist, was auch $x$ sein mag, wenn nur $x$ von den $\alpha$ unabhängig ist, eine symmetrische Function von $\alpha_1,\alpha_2,\ldots,\alpha_n$. Wenn wir also die Multiplication der einzelnen Factoren ausführen und
%CONT>>
